\documentclass{apcalc}
\title{Unit 2 Practice Set: Differentiation: Definition and Fundamental Properties}
\author{Paul Wang}
\date{November 2026}
\volume{Practice Set 2 of 10 (topics 2.1–2.10)}
\shorttitle{Unit 2 --- Practice Set}
\begin{document}
\maketitle
\begin{abstract}
Twelve exam-style multiple-choice questions on Unit 2, grouped by the sections of the Unit 2 Study Guide, and one nine-point free-response question that estimates a derivative from a table and tests differentiability of a piecewise function. All items are no-calculator except those marked [calculator]. Write each answer in the blank beside the number; the key on the last pages gives an answers-only strip, a reason and rule for every item, and the free-response scoring guide.
\end{abstract}
\scoringline{12 MC \quad\blank\ / 9 FR}{Exam pace: about 2 minutes per multiple-choice item (24 min) and 15 minutes for the free response.}
\section*{1\quad Rates of change and the definition of the derivative}
\begin{setrule}R2.1–R2.3. Average rate $=\frac{f(b)-f(a)}{b-a}$. $f'(a)=\lim_{h\to0}\frac{f(a+h)-f(a)}{h}$: recognize the form, name $f$ and $a$.\end{setrule}
\begin{q}{1}
For $f(x)=x^3-2x$, the average rate of change of $f$ on $[1,3]$ is
\chfour{$10$}{$11$}{$22$}{$25$}
\end{q}
\begin{q}{2}
$\displaystyle\lim_{h\to0}\frac{\sqrt{9+h}-3}{h}$ is
\chfour{$0$}{$\dfrac{1}{18}$}{$\dfrac16$}{$\dfrac13$}
\end{q}
\begin{q}{3}
$\displaystyle\lim_{x\to1}\frac{e^{2x}-e^2}{x-1}$ is
\chfour{$e^2$}{$2e$}{$2e^2$}{nonexistent}
\end{q}
\section*{2\quad Estimating derivatives and differentiability}
\begin{setrule}R2.4–R2.6. Estimate with a difference quotient. Differentiable implies continuous; corners, cusps, vertical tangents and jumps are not differentiable.\end{setrule}
\begin{q}{4}
The table gives selected values of a differentiable function $W$. Using the data, the best estimate of $W'(6)$ is
\begin{center}\small\renewcommand{\arraystretch}{1.35}\begin{tabular}{|c|c|c|c|c|}
\hline
$t$ & 0 & 4 & 8 & 11 \\ \hline
$W(t)$ & 1 & 13 & 29 & 23 \\ \hline
\end{tabular}\end{center}
\chfour{$-2$}{$2$}{$3$}{$4$}
\end{q}
\begin{q}{5}
The graph of $f$ is shown. At which of the following values of $x$ is $f$ not differentiable?
\begin{center}\includegraphics[width=0.45\linewidth]{figs/u2-q5.pdf}\end{center}
\chfour{$x=1$ only}{$x=4$ only}{$x=2$ and $x=4$}{$x=2$ only}
\end{q}
\begin{q}{6}
A cup of coffee cools; its temperature $H(t)$, in degrees, is given in the table for $t$ in minutes. Which is the best estimate of $H'(5)$, and its units?\calcmark
\begin{center}\small\renewcommand{\arraystretch}{1.35}\begin{tabular}{|c|c|c|c|c|}
\hline
$t$ (min) & 0 & 3 & 7 & 10 \\ \hline
$H(t)$ (°F) & 180 & 172 & 164 & 160 \\ \hline
\end{tabular}\end{center}
\chtwo{$-\frac83$ degrees per minute}{$-2$ degrees per minute}{$-1$ degree per minute}{$168$ degrees}
\end{q}
\begin{q}{7}
Let $f(x)=\begin{cases}ax^2+b,&x\le2\\ x^3-3,&x>2\end{cases}$. For which $a$ and $b$ is $f$ differentiable at $x=2$?
\chfour{$a=3$, $b=-7$}{$a=3$, $b=17$}{$a=\frac32$, $b=-1$}{$a=12$, $b=-43$}
\end{q}
\section*{3\quad Basic derivative rules}
\begin{setrule}R2.7–R2.9. $\frac{d}{dx}x^n=nx^{n-1}$; $\sin\to\cos$, $\cos\to-\sin$, $e^x\to e^x$, $\ln x\to\frac1x$. Tangent: $y=f(a)+f'(a)(x-a)$.\end{setrule}
\begin{q}{8}
If $y=\dfrac{4}{\sqrt{x}}+x$, then $\dfrac{dy}{dx}$ at $x=4$ is
\chfour{$-\dfrac14$}{$0$}{$\dfrac34$}{$\dfrac54$}
\end{q}
\begin{q}{9}
An equation of the line tangent to $y=3\ln x+x^2$ at $x=1$ is
\chfour{$y=5x-4$}{$y=5x-5$}{$y=2x-1$}{$y=5x-1$}
\end{q}
\section*{4\quad Product and quotient rules}
\begin{setrule}R2.10–R2.12. $(uv)'=u'v+uv'$; $(u/v)'=\frac{u'v-uv'}{v^2}$; $\tan\to\sec^2$, $\sec\to\sec\tan$.\end{setrule}
\begin{q}{10}
If $g(x)=\dfrac{3x-1}{x^2+1}$, then $g'(1)=$
\chfour{$-\dfrac12$}{$\dfrac12$}{$\dfrac32$}{$\dfrac52$}
\end{q}
\begin{q}{11}
If $f(x)=x^2\sin x$, then $f'\!\left(\frac{\pi}{2}\right)=$
\chfour{$0$}{$\dfrac{\pi}{2}$}{$\dfrac{\pi^2}{4}$}{$\pi$}
\end{q}
\begin{q}{12}
If $y=\tan x\sec x$, then $\dfrac{dy}{dx}$ at $x=\frac{\pi}{4}$ is
\chfour{$\sqrt2$}{$2\sqrt2$}{$3\sqrt2$}{$2\sqrt2+2$}
\end{q}
\clearpage\section*{Free response}
Show all work. Justify every conclusion with calculus, and give decimal answers correct to three places.
\begin{frq}{Question 1 --- Rates from a table; a piecewise model (9 points, no calculator)}
Water flows into a tank at a rate of $R(t)$ liters per hour, where $R$ is differentiable and $t$ is measured in hours. Selected values of $R$ are given in the table.
\partlabel{a}Find the average rate of change of $R$ over $0\le t\le10$. Indicate units.\workspace{1.6in}
\partlabel{b}Use the data to estimate $R'(6.5)$. Show the computations that lead to your answer, and explain the meaning of your estimate in context.\workspace{1.6in}
\partlabel{c}For $10\le t\le14$ the rate is modeled by $R(t)=\begin{cases}48-a(t-10)^2,&10\le t\le12\\ bt+c,&12<t\le14\end{cases}$ where $R$ is differentiable at $t=12$ and $R(14)=24$. Find the values of $a$, $b$ and $c$.\workspace{1.6in}
\end{frq}
\clearpage\section*{Answer key}
\keystrip{\keyitem{1}{B}\keyitem{2}{C}\keyitem{3}{C}\keyitem{4}{D}\keyitem{5}{D}\keyitem{6}{B}\keyitem{7}{A}\keyitem{8}{C}\keyitem{9}{A}\keyitem{10}{B}\keyitem{11}{D}\keyitem{12}{C}}
\vspace{4pt}
\begin{keytable}
\keysection{1. Rates of change and the definition of the derivative}
1 & B & $\frac{f(3)-f(1)}{3-1}=\frac{21-(-1)}{2}=11$. & R2.1\\
2 & C & It is $f'(9)$ for $f(x)=\sqrt x$: $\frac{1}{2\sqrt9}=\frac16$. & R2.3\\
3 & C & It is $f'(1)$ for $f(x)=e^{2x}$: $2e^{2}$. & R2.3\\
\keysection{2. Estimating derivatives and differentiability}
4 & D & $\frac{W(8)-W(4)}{8-4}=\frac{29-13}{4}=4$. & R2.4\\
5 & D & Corner at $x=2$ (slopes 1 and 0); at $x=4$ the slopes are 0 on both sides. & R2.5\\
6 & B & $\frac{H(7)-H(3)}{7-3}=\frac{164-172}{4}=-2$ °/min. & R2.4, \emph{calc}\\
7 & A & Slopes: $4a=12$, so $a=3$; values: $4a+b=5$, so $b=-7$. & R2.6\\
\keysection{3. Basic derivative rules}
8 & C & $y'=-2x^{-3/2}+1=-\frac14+1=\frac34$. & R2.7\\
9 & A & Point $(1,1)$, slope $\frac31+2=5$: $y=1+5(x-1)$. & R2.9\\
\keysection{4. Product and quotient rules}
10 & B & $\frac{3(2)-(2)(2)}{2^2}=\frac{2}{4}=\frac12$. & R2.11\\
11 & D & $f'=2x\sin x+x^2\cos x$; at $\frac\pi2$: $\pi\cdot1+0=\pi$. & R2.10\\
12 & C & $\sec^3x+\tan^2x\sec x=2\sqrt2+\sqrt2=3\sqrt2$. & R2.12\\
\end{keytable}
\section*{Free-response scoring guide}
\subsection*{Question 1: Rates from a table; a piecewise model (9 points)}
\textbf{(a)} $\frac{48-40}{10}=0.8$ liters per hour per hour.
\begin{itemize}[leftmargin=1.6em,itemsep=0pt]
\item 1 point: $\frac{R(10)-R(0)}{10-0}=0.8$
\item 1 point: units: liters per hour per hour
\end{itemize}
\textbf{(b)} $R'(6.5)\approx\frac{R(8)-R(5)}{8-5}=\frac{55-61}{3}=-2$. At $t=6.5$ hours, the rate at which water flows in is decreasing at about 2 liters per hour per hour.
\begin{itemize}[leftmargin=1.6em,itemsep=0pt]
\item 1 point: difference quotient on $[5,8]$
\item 1 point: estimate $-2$
\item 1 point: interpretation: at $t=6.5$ the inflow rate is decreasing by about 2 liters per hour per hour
\end{itemize}
\textbf{(c)} Continuity at 12: $48-4a=12b+c$. Matching slopes: $-2a(12-10)=-4a=b$. $R(14)=14b+c=24$. Subtract the continuity equation: $2b=24-(48-4a)=4a-24$. With $b=-4a$: $-8a=4a-24$, so $a=2$. Then $b=-8$ and $c=24-14b=136$. Check: $R(12)=48-8=40=12(-8)+136$.
\begin{itemize}[leftmargin=1.6em,itemsep=0pt]
\item 1 point: continuity equation $48-4a=12b+c$
\item 1 point: slope equation $-4a=b$
\item 1 point: uses $R(14)=24$: $14b+c=24$
\item 1 point: solves: $a=2$, $b=-8$, $c=136$
\end{itemize}
\end{document}